% Part: first-order-logic
% Chapter: model-theory
% Section: partial-iso

\documentclass[../../../include/open-logic-section]{subfiles}

\begin{document}

% SEGMENT OLP-0190-S01

\olfileid{mod}{bas}{dlo}
\section{செறிந்த நேரியல் வரிசைகள்}

\begin{defn}
  \emph{முனைகளற்ற செறிந்த நேரியல் வரிசை} என்பது ஒரே ஒரு 2-இட
  !!{predicate}~$<$ ஐக் கொண்ட !!{language} க்கான, பின்வரும்
  வாக்கியங்களை நிறைவுறுத்தும் !!a{structure} $\Struct{M}$:
  \begin{enumerate}
  \item $\lforall[x][\lnot x < x]$;
  \item $\lforall[x][\lforall[y][\lforall[z][(x < y \lif (y < z \lif x
    <z ))]]]$;
  \item $\lforall[x][\lforall[y][(x< y \lor \eq[x][y] \lor y < x)]]$;
  \item $\lforall[x][\lexists[y][x < y]]$;
  \item $\lforall[x][\lexists[y][y < x]]$;
  \item $\lforall[x][\lforall[y][(x < y \lif \lexists[z][(x < z \land
        z < y))]]]$.
 \end{enumerate}
\end{defn}

\begin{thm}\ollabel{thm:cantorQ}
  எந்த இரு !!{enumerable} முனைகளற்ற செறிந்த நேரியல் வரிசைகளும்
  ஓரினமானவை.
\end{thm}

\begin{proof}
  $\Struct{M_1}$ மற்றும் $\Struct{M_2}$ ஆகியவை !!{enumerable}
  முனைகளற்ற செறிந்த நேரியல் வரிசைகள் ஆகட்டும்; ${<_1} =
  \Assign{<}{M_1}$ மற்றும் ${<_2} =
  \Assign{<}{M_2}$ என்றும், அவற்றுக்கிடையிலான எல்லாப் பகுதி
  ஓரினச்சார்புகளின் கணம் $\PIso{I}$ என்றும் கொள்க. குறைந்தபட்சம்
  $\emptyset \in \PIso{I}$ என்பதால் $\PIso{I}$ வெற்றல்ல.
  $\PIso{I}$ முன்னும் பின்னும் பண்பை நிறைவுறுத்துகிறது என்று காட்டுவோம்.
  அப்போது $\Struct{M_1} \iso[p] \Struct{M_2}$; எனவே
  \olref[pis]{thm:p-isom1} படி தேற்றம் பின்வரும்.

  $\PIso{I}$ முன்னும் பண்பை நிறைவுறுத்துகிறது என்று காட்ட,
  $p \in
  \PIso{I}$ என்றும், $i = 1$, \dots,~$n$ க்கு $p(a_i) = b_i$ என்றும்
  கொள்க; பொதுமை இழப்பின்றி $a_1 <_1 a_2 <_1 \cdots <_1
  a_n$ என்று கொள்க. $a \in \Domain{M_1}$ கொடுக்கப்பட்டால்,
  $b \in \Domain{M_2}$ ஐப் பின்வருமாறு கண்டறிக:
  \begin{enumerate}
  \item $a <_1 a_1$ எனில், $b <_2
    b_1$ ஆகுமாறு $b \in \Domain{M_2}$ ஐக் கொள்க;
  \item $a_n <_1 a$ எனில், $b_n <_2 b$ ஆகுமாறு
    $b \in \Domain{M_2}$ ஐக் கொள்க;
 \item ஏதேனும் $i$ க்கு $a_i <_1 a <_1 a_{i+1}$ எனில்,
   $b_i <_2 b <_2 b_{i+1}$ ஆகுமாறு $b \in
   \Domain{M_2}$ ஐக் கொள்க.
  \end{enumerate}
  $\Struct{M_2}$ ஒரு முனைகளற்ற செறிந்த நேரியல் வரிசை என்பதால்,
  வேண்டிய பண்புடைய $b$ ஒன்றை எப்போதும் கண்டறியலாம்.
  $q \in \PIso{I}$ ஆனது $p$ இன் வேண்டிய நீட்சியாக இருக்குமாறு
  $q = p \cup \{ \langle a, b \rangle \}$ என்று வரையறுக்க.
  இதனால் முன்னும் பண்பு நிறுவப்படுகிறது. பின்னும் பண்பு இதே போன்றது.
  எனவே $\Struct{M_1} \iso[p] \Struct{M_2}$; மேலும்
  \olref[pis]{thm:p-isom1} படி, $\Struct{M_1} \iso
  \Struct{M_2}$.
\end{proof}

\begin{prob}
  $\PIso{I}$ பின்னும் பண்பை நிறைவுறுத்துகிறது என்று சரிபார்த்து,
  \olref[mod][bas][dlo]{thm:cantorQ} இன் நிரூபிப்பை நிறைவுசெய்க.
\end{prob}

\begin{rem}
  $\Struct{S}$ ஏதேனும் ஓர் !!{enumerable} முனைகளற்ற செறிந்த நேரியல்
  வரிசையாகட்டும். அப்போது (\olref{thm:cantorQ} படி)
  $\Struct{S} \iso
  \Struct{Q}$; இங்கே $\Struct{Q} = (\Rat, <)$ என்பது விகிதமுறு எண்களின்
  கணம் $\Rat$ ஐ ஆட்களமாகக் கொண்ட !!{enumerable} செறிந்த நேரியல் வரிசை.
  இப்போது \olref[thm]{remark:R} இலிருந்து !!{structure}
  ~$\Struct{R} =
  (\Real, <)$ ஐ மீண்டும் கருதுக. $\Struct{R}
  \elemequiv \Struct{S}$ ஆகுமாறு !!a{enumerable} !!{structure}
  ~$\Struct{S}$ ஒன்று உள்ளது என்று கண்டோம். ஆனால் $\Struct{S}$ ஒரு
  !!a{enumerable} முனைகளற்ற செறிந்த நேரியல் வரிசை; எனவே அது
  !!{structure}~$\Struct{Q}$ க்கு ஓரினமானது (ஆகவே அடிப்படையாகச்
  சமானமானது). அடிப்படைச் சமானத்தின் கடப்புநிலைப் பண்பின்படி,
  $\Struct{R} \elemequiv
  \Struct{Q}$. (அதே முன்னும் பின்னும் வாதத்தால்
  $\Struct{R} \iso[p] \Struct{Q}$ என்பதை நிறுவி, இதை நேரடியாகவும்
  காட்டியிருக்கலாம்.)
\end{rem}
\end{document}
